\section{The test} \label{sec:test}

\subsection{Data and processing}

The primary is the first released SPOC 2-min light curve, column \code{PDCSAP\_FLUX} (corrected flux), quality mask
17087. Only if it is absent on 2027 February 28 may the primary use QLP 200-s (mask 1610619071), then TESS-SPOC
200-s (\code{PDCSAP\_FLUX}, mask 17087) if QLP is also absent, using the version public on that date and a saved
MAST search record (Appendix~\ref{app:data}). QLP uses \code{SYS\_RM\_FLUX}, or \code{DET\_FLUX} only if the former
column is absent (Appendix~\ref{app:commands}). An existing primary with no counted transit gives O4; it never
triggers a product switch.

The biweight detrending filter has window $W=0.75$\,d \citep{Hippke2019}; both signals are masked over
$\pm1.5\,\Tfourteen$. Upward outliers above $4\times1.4826$ times the median absolute deviation are removed.
Segments split at gaps longer than 0.3\,d in the cleaned series with both transit masks removed. The first and last
$W/2$ (9\,h) of each segment are edge zones. A transit is counted only if (i) at least 80\,\% of its window has
cadences after cleaning, (ii) at least 2 cadences lie on each side, 0.75 to 3\,$\Tfourteen$ from its predicted
mid-time, (iii) its whole window lies within one segment outside the edge zones, and (iv) it is a frozen Sector 110
epoch. Otherwise it is ``not tested'' and never counts against the signal (Appendix~\ref{app:processing}).

\subsection{Statistic and thresholds} \label{sec:stat}

The depth of one transit is 1 minus the mean detrended flux inside $|t - t_c| < \Tfourteen/2$, where $t_c$ is its
predicted mid-time. The test statistic $d$ is the mean of these depths over the $n$ counted transits. Nothing is
searched in period or in time.

The noise is measured on the same light curve. The same window is placed at fake transit times between the real ones,
and each fake ephemeris goes through the full detrending, masking and counting; the windows obtained form the noise
pool (Appendix~\ref{app:noise}). The pool is resampled 20\,000 times for each of 20 seeds; the medians of the resulting thresholds and sensitivity
are used (Appendix~\ref{app:thresholds}).
From these resamples, for the $n$ actually counted:
\begin{itemize}
\item the detection threshold $\ddet$ is the 99.9th percentile of the mean of $n$ noise windows;
\item the band of depths consistent with the prediction runs from $\dlo$, the 0.1st percentile of noise plus a depth
  drawn from $N(D, D_{\rm err})$, to $\dhi$, the 99.9th percentile of noise plus a depth drawn from
  $N(D_{98}, D_{98,\rm err})$;
\item the sensitivity $\Sens = P(\mathrm{detect} \mid \mathrm{real})$ is the fraction of draws of noise plus
  $N(D, D_{\rm err})$ above $\ddet$.
\end{itemize}
The percentile is not the false-alarm level: because neighbouring noise windows overlap, the bootstrap underestimates
the tails. The false-detection levels were therefore measured on archival data with a separate set of fake ephemerides:
0.47\,\% for A and 0.49\,\% for B overall in SPOC 2-min data, with the level for each number of counted transits in
Table~\ref{tab:fa}. These trials are not independent of each other, so their binomial bounds are approximate.

\subsection{Four outcomes} \label{sec:outcomes}

Each signal receives exactly one outcome (Table~\ref{tab:outcomes}).

\begin{longtable}{>{\raggedright\arraybackslash}p{0.17\textwidth}>{\raggedright\arraybackslash}p{0.43\textwidth}>{\raggedright\arraybackslash}p{0.32\textwidth}}
\caption{Outcome of each signal. All quantities are computed for the $n$ counted transits.\label{tab:outcomes}}\\
\toprule
Outcome & Condition & Meaning \\
\midrule
\endfirsthead
\toprule
Outcome & Condition & Meaning \\
\midrule
\endhead
\textbf{O1 recovered} & $d > \ddet$, $\dlo \leq d \leq \dhi$, and the timing is consistent (Section~\ref{sec:timing}) &
  a dip consistent with the prediction is present \\
\textbf{O2 inconsistent} & $d > \ddet$, and at least one of: $d < \dlo$; $d > \dhi$; the timing is inconsistent & a dip is present, but
  not the predicted one (the failed condition is recorded) \\
\textbf{O3 excluded} & $d \leq \ddet$, $d < \dlo$, and $\Sens \geq 0.95$ & the predicted signal is not present at the
  predicted depth, and the data had at least 95\,\% sensitivity to it \\
\textbf{O4 insufficient} & anything else: fewer than 10 usable cadences, or no registered QLP flux column
  (Appendix~\ref{app:commands}); $n = 0$; a noise pool of fewer than 100 windows; or $d \leq \ddet$ and either $d \geq \dlo$ or
  $\Sens < 0.95$ & the data cannot decide \\
\bottomrule
\end{longtable}

Transits that are not tested reduce $n$; they never count against the signal. O1 and O2 do not depend on the
sensitivity: a $d$ above $\ddet$ gives O1 or O2 whatever the sensitivity; its false-detection level is the one in
Table~\ref{tab:fa} for that $n$. The sensitivity matters only when $d \leq \ddet$. In the dry run on the archival
Sector 32 data the sensitivity for A was 0.92, below the 0.95 needed for O3 (Appendix~\ref{app:validation}). Even a
clean non-detection of A may therefore give O4. For B, with two or three transits, the sensitivity may be far below
0.95. In that case an O1 of B is still possible but an O3 is not, and B will most likely end as O4 unless Sector 110 is
quieter than Sector 32.

O4 means only that the data cannot decide; it says nothing against the signal. O1 means that a dip consistent with the prediction is
present in new data; it does not show that the signal comes from a planet (Section~\ref{sec:limits}).

\subsection{Timing rule} \label{sec:timing}

The timing is checked only when $d > \ddet$.
\begin{itemize}
\item The common time shift $\Delta t$ of all counted transits is the shift between $-30$ and $+30$\,min, in steps of
  0.5\,min, that maximises their mean depth. If several shifts share the maximum (a plateau, because the grid is finer
  than the cadence), the middle one is taken; for an even number of them, the later of the two middle shifts.
\item $\sigma_{\rm meas}$ is the standard deviation of the recovered shift over 300 injections of a box of depth $d$ into
  $n$ random noise-pool windows, with a true shift of 0 and a fixed seed.
\item $\sigma_{\rm pred}$ is the mean of the prediction uncertainties (Tables~\ref{tab:timesA} and \ref{tab:timesB}) of
  the counted epochs.
\item The timing is consistent if $|\Delta t| \leq 3\sqrt{\sigma_{\rm pred}^2 + \sigma_{\rm meas}^2}$. For A this
  tolerance is at least 12.3\,min.
\item The shift of each counted transit alone is reported. It does not enter the outcome.
\end{itemize}

\subsection{Processing variants}

Up to six fixed processing variants are run after the primary analysis and reported next to it: (a) and (b)
detrending windows of 0.5 and 1.0\,d; (c) cadences near SPOC desaturation flags removed within $\pm1$\,h (SPOC primary only); (d) the QLP
light curve (not for a QLP primary; possibly run later, on the release public on 2027 February 28); (e) the uncorrected
simple-aperture flux; (f) the counting rule of an earlier draft, for comparison only. They never
change the primary outcome. They only decide whether an O1 or an O3 is called ``robust'', a word with a narrow meaning:
among the variants (a) to (e) that were run and reach $\Sens \geq 0.95$, none returned O3 against a primary O1, and
none returned O1 or O2 against a primary O3. If none of them reaches 0.95, the qualifier is ``robustness not
testable'', and while variant (d) is pending, ``robust'' and ``robustness not testable'' remain provisional;
a qualifying contradiction already makes the result ``not robust''. The qualifier ``robust'' does not mean that all variants agree
(Appendix~\ref{app:variants}).
